\section{Experiments}
\label{sec:results}

\subsection{Protocol}
\label{sec:protocol}

\paragraph{Set-up.} All GPU experiments use the official GainLoRA code on SuperNI
\citep{wang2022superni}; a launcher sets the switches of Eq.~\eqref{eq:gates} and leaves the published
methods InfLoRA \citep{liang2024inflora} and GainLoRA \citep{liang2025gainlora} unchanged. We use T5-base
\citep{raffel2020t5} and 10 epochs per task instead of T5-large and 100, to fit free GPUs; a few runs use 100.
\textbf{Because of the reduced setting, none of our numbers can be compared with the tables of the
published papers.} Runs without definitions replace the definition of every task by the empty string,
for training and for testing. Before each run we fixed its rules and a \emph{line}: what would count
for or against the method.

\paragraph{Sequences and inputs.} The five tasks (relation extraction, SST-2, cause-event detection,
Sentiment140, emotion classification) have 800 training and 100 test examples each; a 15-task run (one
seed) is in Section~\ref{sec:fifteen} of the appendix. \textbf{They were not drawn at random}: the first
three are short-input tasks that InfLoRA had forgotten most. The unrelated inputs $\mathcal{B}$ are
4{,}000 inputs from ten tasks of another benchmark (Section~\ref{sec:background}). Each task has three
reworded definitions (\emph{light}, \emph{medium}, \emph{heavy}), \textbf{written by machine during
development; two design choices were made while looking at where they were routed}
(Appendix~\ref{app:history}). They are a development set, not an independent test.

\paragraph{Metrics.} AP is the average Rouge-L after the last task, FT the forgetting, and Learned the
score right after each task was learned, so that forgetting and little learning can be told apart.
\emph{Answers changed} counts test inputs of earlier tasks whose final answer differs, as a string,
from the answer right after their task was learned. Besides the official test (one random reference)
we report a \emph{re-test} against all references (Appendix~\ref{app:cut}).

\subsection{Without instructions}
\label{sec:noinst}

\paragraph{The published methods without their definitions.} We run InfLoRA and GainLoRA again on the
five tasks without definitions (three seeds, Table~\ref{tab:noinstseedsshort}). Both hold their answers
less well, by a line fixed before the run: they change \AudSeedsGainloraChanged{} (GainLoRA) and
\AudSeedsInfloraChanged{} (InfLoRA) of \AudSeedsCompared{} answers, against \FiveGainChanged{} and
\FiveInfChanged{} with definitions, and forgetting is \AudSeedsGainloraFT{} and \AudSeedsInfloraFT{}
points against \FiveGainFT{} and \FiveInfFT{}. No benchmark input occurs in two tasks
(\AudConfOneShared{} of \AudConfOneInputs{} inputs of the 15 SuperNI tasks; Table~\ref{tab:conflictcount},
appendix), so no loss is forced by shared texts. The line, written before: a method that
learned (a score right after learning of at least 60\% of its score with definitions) ``holds its
answers less well'' if forgetting rises by at least five points, or the changed answers at least double.
This is a first look at two methods with a tenth of their epochs; it says nothing about other methods.

\begin{table}[t]
  \caption{The five tasks without definitions, T5-base, 10 epochs, means over \AudSeedsN{} seeds (changed answers summed), re-test (all references). Answers changed: answers to test inputs of the first four tasks that differ at the end. The 100-epoch rows are one seed.}
  \label{tab:noinstseedsshort}
  \centering
  {\footnotesize\setlength{\tabcolsep}{2.5pt}\begin{tabular}{lrrrr}
    \toprule
    Method & \shortstack[r]{Score after\\learning $\uparrow$} & \shortstack[r]{Final\\score $\uparrow$} & \shortstack[r]{Forgetting\\(points) $\downarrow$} & \shortstack[r]{Answers changed\\of \AudSeedsCompared{} $\downarrow$} \\
    \midrule
    InfLoRA & \AudSeedsInfloraLearned{} & \AudSeedsInfloraAP{} {\scriptsize (\AudSeedsInfloraAPLow{}--\AudSeedsInfloraAPHigh{})} & \AudSeedsInfloraFT{} & \AudSeedsInfloraChanged{} \\
    GainLoRA & \AudSeedsGainloraLearned{} & \AudSeedsGainloraAP{} {\scriptsize (\AudSeedsGainloraAPLow{}--\AudSeedsGainloraAPHigh{})} & \AudSeedsGainloraFT{} & \AudSeedsGainloraChanged{} \\
    InfLoRA, 100 epochs & \AudFullInfloraLearned{} & \AudFullInfloraAP{} & \AudFullInfloraFT{} & \AudFullInfloraChanged{} of \AudFullCompared{} \\
    GainLoRA, 100 epochs & \AudFullGainloraLearned{} & \AudFullGainloraAP{} & \AudFullGainloraFT{} & \AudFullGainloraChanged{} of \AudFullCompared{} \\
    InfLoRA, 100 epochs, with definitions & \AudFullDefOriginalInfloraLearned{} & \AudFullDefOriginalInfloraAP{} & \AudFullDefOriginalInfloraFT{} & \AudFullDefOriginalInfloraChanged{} of \AudFullDefCompared{} \\
    GainLoRA, 100 epochs, with definitions & \AudFullDefOriginalGainloraLearned{} & \AudFullDefOriginalGainloraAP{} & \AudFullDefOriginalGainloraFT{} & \AudFullDefOriginalGainloraChanged{} of \AudFullDefCompared{} \\
    \midrule
    \plain{} & \AudSeedsPlainLearned{} & \AudSeedsPlainAP{} {\scriptsize (\AudSeedsPlainAPLow{}--\AudSeedsPlainAPHigh{})} & \AudSeedsPlainFT{} & \AudSeedsPlainChanged{} \\
    \centre{} & \AudSeedsCentreLearned{} & \AudSeedsCentreAP{} {\scriptsize (\AudSeedsCentreAPLow{}--\AudSeedsCentreAPHigh{})} & \AudSeedsCentreFT{} & \AudSeedsCentreChanged{} \\
    \recog{} (ours) & \AudSeedsRidgeLearned{} & \AudSeedsRidgeAP{} {\scriptsize (\AudSeedsRidgeAPLow{}--\AudSeedsRidgeAPHigh{})} & \AudSeedsRidgeFT{} & \AudSeedsRidgeChanged{} \\
    \handicap{} (ours) & \AudSeedsHandicapLearned{} & \AudSeedsHandicapAP{} {\scriptsize (\AudSeedsHandicapAPLow{}--\AudSeedsHandicapAPHigh{})} & \AudSeedsHandicapFT{} & \AudSeedsHandicapChanged{} \\
    \recog{} (ours), 100 epochs & \AudFullOursRidgeLearned{} & \AudFullOursRidgeAP{} & \AudFullOursRidgeFT{} & \AudFullOursRidgeChanged{} of \AudFullCompared{} \\
    \handicap{} (ours), 100 epochs & \AudFullOursHandicapLearned{} & \AudFullOursHandicapAP{} & \AudFullOursHandicapFT{} & \AudFullOursHandicapChanged{} of \AudFullCompared{} \\
    own task (ceiling) & \AudSeedsOwnLearned{} & \AudSeedsOwnAP{} {\scriptsize (\AudSeedsOwnAPLow{}--\AudSeedsOwnAPHigh{})} & \AudSeedsOwnFT{} & \AudSeedsOwnChanged{} \\
    \bottomrule
\end{tabular}}

\end{table}

\paragraph{Recognition.} The feature matters: with the zones built from inputs
(Section~\ref{sec:boundmain}) on the 15 SuperNI tasks, frozen
T5-base states give \AudZoneFrozenTOwn{}\% of the test inputs to their own task and block
\AudZoneFrozenTBlocked{}\%, a sentence encoder \AudZoneSentenceOwn{}\% and \AudZoneSentenceBlocked{}\%.
Table~\ref{tab:recognisershort} gives the routes of \recog{} and \handicap{}: the
first row is development data on which the rule was chosen; the Long Sequence and the 30 unseen
tasks were read once against a line fixed before; the 120 tasks include the first row.
On the 15 tasks the audit says \AudDevLLamOneAuditEst{}\% for the average task
(95\% bound \AudDevLLamOneAuditBound{}\%), next to the \AudDevLLamOneChanged{}\% that changed. The floor at zero of \handicap{} was added after seeing numbers on
these 15 tasks, so its numbers there are development numbers. Two
tasks built on the same texts split their inputs under \recog{} (\AudScalePairOlderLEnd{}\% and
\AudScalePairNewerLEnd{}\% reach their own task); under \handicap{} the older keeps
\AudScalePairOlderHEnd{}\% and the newer gets \AudScalePairNewerHEnd{}\%. \AudConfPairShare{}\% of
their texts are the same word for word, and there no rule without instructions can do better.

\begin{table}[t]
  \caption{Routes of \recog{} and \handicap{} (\%, means over tasks): test inputs at their own task after their task arrived and after the last task; routes changed; the audit's 95\% bound. Routes only; the first row is development data.}
  \label{tab:recognisershort}
  \centering
  {\footnotesize\setlength{\tabcolsep}{1.2pt}\begin{tabular}{lrrrrrr}
    \toprule
    & \multicolumn{4}{c}{\recog{}, \%} & \multicolumn{2}{c}{\handicap{}, $\nu = 2$, \%} \\
    \cmidrule(lr){2-5}\cmidrule(l){6-7}
    & \shortstack[r]{own task\\at arrival $\uparrow$} & \shortstack[r]{own task\\at the end $\uparrow$} & \shortstack[r]{routes\\changed $\downarrow$} & \shortstack[r]{audit\\95\% bound} & \shortstack[r]{own task\\at the end $\uparrow$} & \shortstack[r]{routes\\changed $\downarrow$} \\
    \midrule
    15 SuperNI tasks (development) & \AudDevLLamOneArrive{} & \AudDevLLamOneEnd{} & \AudDevLLamOneChanged{} & \AudDevLLamOneAuditBound{} & \AudDevHPrimeTTwoEnd{} & \AudDevHPrimeTTwoChanged{} \\
    15 Long Sequence tasks & \AudScaleLongThreeArrive{} & \AudScaleLongThreeEnd{} & \AudScaleLongThreeChanged{} & \AudScaleLongThreeAuditBound{} & \AudScaleLongThreeHEnd{} & \AudScaleLongThreeHChanged{} \\
    30 unseen tasks (read once) & \AudFreshFourthLArrive{} & \AudFreshFourthLEnd{} & \AudFreshFourthLChanged{} & \AudFreshFourthAuditBound{} & \AudFreshFourthHEnd{} & \AudFreshFourthHChanged{} \\
    90 fresh tasks & \AudScaleFreshNinetyArrive{} & \AudScaleFreshNinetyEnd{} & \AudScaleFreshNinetyChanged{} & \AudScaleFreshNinetyAuditBound{} & \AudScaleFreshNinetyHEnd{} & \AudScaleFreshNinetyHChanged{} \\
    120 tasks & \AudScaleAllHundredTwentyArrive{} & \AudScaleAllHundredTwentyEnd{} & \AudScaleAllHundredTwentyChanged{} & \AudScaleAllHundredTwentyAuditBound{} & \AudScaleAllHundredTwentyHEnd{} & \AudScaleAllHundredTwentyHChanged{} \\
    \bottomrule
\end{tabular}}

\end{table}

\paragraph{Trained answers.} Adapters trained as usual, behind \recog{} and \handicap{} ($\nu = 2$), were run on the
five tasks without definitions, with routes written down before the run. The line: the final score at
least two points above the better published method, forgetting below two points, and at most a
quarter of the answers that GainLoRA changes. \recog{} and \handicap{} change \AudSeedsRidgeChanged{} and
\AudSeedsHandicapChanged{} of \AudSeedsCompared{} answers and end above both published methods in
\AudSeedsRidgeAboveInAll{} of \AudSeedsN{} seeds. By this line fixed before the run, read on the re-test (the official script
cannot apply routes), they beat the published methods in this setting. The audit estimated
\AudRouteFiveRidgeFirstFourAuditEst{}\% changed routes; \AudRouteFiveRidgeChangedShare{}\% changed,
within the 95\% bound of each task. \recog{} and \handicap{} end at \AudSeedsRidgeAPOne{} and
\AudSeedsHandicapAPOne{}, with forgetting of \AudSeedsRidgeFT{} and \AudSeedsHandicapFT{} points. Caveats: at 10 epochs most of that gap is learning; even \centre{} ends at \AudSeedsCentreAPOne{}; it is one
sequence of five development tasks. On fifteen tasks (one seed, our rules only;
Section~\ref{sec:nodeftrained}) \recog{} ends \AudFifteenNoDefRidgeBelowOwn{} points below the ceiling, with
\AudFifteenNoDefRidgeChanged{} of \AudFifteenNoDefCompared{} changed answers.

\paragraph{The published methods with their published epochs.} With \AudFullEpochs{} epochs per
task (seed 42; lines fixed before the run) the two methods learn the five tasks without definitions:
\AudFullInfloraLearnedOne{} and \AudFullGainloraLearnedOne{} right after learning, against
\AudOursRidgeLearnedOne{} for ours with 10. Then they lose them: they forget \AudFullInfloraFTInt{}
and \AudFullGainloraFTInt{} points and end at \AudFullInfloraAPOne{} and \AudFullGainloraAPOne{}, against
\AudFullOursRidgeAPOne{} for \recog{} with the same \AudFullEpochs{} epochs (Table~\ref{tab:noinstseedsshort}; Figures~\ref{fig:answergrid}--\ref{fig:learnkeep}, Appendix~\ref{app:pictures}). By a line fixed before, \recog{} beats them at equal epochs.
Here the gap is forgetting. With their definitions left in, the same methods at \AudFullEpochs{} epochs forget \AudFullDefOriginalInfloraFTOne{} and
\AudFullDefOriginalGainloraFTOne{} points and change \AudFullDefOriginalInfloraChanged{} and \AudFullDefOriginalGainloraChanged{} answers: the instruction was what held their answers.

\subsection{With instructions: the certified switch}
\label{sec:withinstr}

\begin{table}[!htbp]
  \caption{Five tasks, original instructions (T5-base, 10 epochs; re-test). Answers changed: test inputs of the first four tasks whose
  final answer differs from the answer right after learning; certified inputs: inputs inside a zone. Means over
  \FiveCertSeeds{} seeds (lowest and highest in brackets); counts summed over seeds.}
  \label{tab:fivemain}
  \centering
  \footnotesize
  \setlength{\tabcolsep}{3pt}
  \begin{tabular}{lccccc}
    \toprule
    Method & \shortstack[c]{Score\\right after\\learning $\uparrow$} & \shortstack[c]{Final\\score $\uparrow$} & \shortstack[c]{Forgetting\\(points) $\downarrow$} & \shortstack[c]{Answers\\changed $\downarrow$} & Certified inputs \\
    \midrule
    \cert{} (ours) & 79.74 {\scriptsize(79.2--80.8)} & 79.74 {\scriptsize(79.2--80.8)} & 0.00 & 0 of 1{,}200 & 1{,}500 of 1{,}500 \\
    \nearest{} & 79.74 {\scriptsize(79.2--80.8)} & 79.74 {\scriptsize(79.2--80.8)} & 0.00 & 0 of 1{,}200 &  \\
    \centre{} & 79.74 {\scriptsize(79.2--80.8)} & 79.74 {\scriptsize(79.2--80.8)} & 0.00 & 0 of 1{,}200 &  \\
    \keyswitch{} & 79.74 {\scriptsize(79.2--80.8)} & 79.74 {\scriptsize(79.2--80.8)} & 0.00 & 0 of 1{,}200 &  \\
    \plain{} & 79.74 {\scriptsize(79.2--80.8)} & 36.08 {\scriptsize(32.3--39.5)} & 54.58 {\scriptsize(51.6--58.6)} & 994 of 1{,}200 &  \\
    \midrule
    InfLoRA & 60.02 {\scriptsize(56.6--63.3)} & 62.93 {\scriptsize(59.0--65.4)} & $-$3.14 {\scriptsize($-$5.3 to $-$1.7)} & 370 of 1{,}200 &  \\
    GainLoRA & 61.42 {\scriptsize(59.1--64.4)} & 61.96 {\scriptsize(60.9--63.9)} & $-$0.51 {\scriptsize($-$2.6 to 0.7)} & 65 of 1{,}200 &  \\
    \bottomrule
\end{tabular}

\end{table}

\paragraph{Five tasks.}\label{sec:five} The certified switch keeps every answer: after the fifth task the average score
equals the score right after learning, \FiveCertAP{} Rouge-L, forgetting is \FiveCertFT{}, and none of
the \FiveCertCompared{} answers to the test inputs of the first four tasks differs from the answer
given right after the task was learned (Table~\ref{tab:fivemain}). All \FiveCertInputsOne{} test inputs
lie inside their own zone (Proposition~\ref{prop:own}), and no training input lay inside an older zone
(\FiveCertBlocked{} of \FiveCertTrainInputs{}). This follows from Theorem~\ref{thm:writeonce}; the
experiment confirms the implementation. \nearest{}, \centre{} and \keyswitch{} give the same numbers.
The same adapters with all switches on (\plain{}) fall from \FivePlainLearned{} to \FivePlainAP{}, and
\FivePlainChanged{} of \FivePlainCompared{} answers change. InfLoRA and GainLoRA change \FiveInfChanged{}
and \FiveGainChanged{} of \FiveInfCompared{} answers, with forgetting near zero (\FiveInfFT{} and
\FiveGainFT{}); they end at \FiveInfAP{} and \FiveGainAP{} because they start lower
(\FiveInfLearned{} and \FiveGainLearned{} right after learning, against \FiveCertLearned{}).
\textbf{With 10 epochs they are at the level of guessing on the two two-label tasks, so this comparison
cannot show how well they retain what they learn} (Appendix~\ref{app:impl}). On 15 tasks (one seed; Section~\ref{sec:fifteen}) the switch changed
\FifteenCertChanged{} of \FifteenCertCompared{} answers (AP \FifteenCertAP{}, plain adapters
\FifteenPlainAP{}); the published methods learned less and forgot \FifteenInfFT{} and \FifteenGainFT{} points.

\paragraph{Reworded instructions.}\label{sec:reworded} A benchmark input carries its definition word for
word, so any rule that recognises it looks perfect; the test is to ask the same task in other words
(Appendix~\ref{sec:fivefull}). No answer changes under any wording. With the light rewording this is
guaranteed (all \FiveCertZoneOneLight{} inputs are still certified). With the medium and heavy rewording it
is \emph{not}: only \FiveCertZoneOneMedium{} and \FiveCertZoneOneHeavy{} of the \FiveCertInputsOne{}
inputs are inside a zone. The others kept their answers because no later task came nearer
(Proposition~\ref{prop:uncertified}); in another order of tasks 31 routes of heavily reworded inputs
would change (Section~\ref{sec:orders}). The average falls from \FiveCertAP{} to \FiveCertAPLight{},
\FiveCertAPMedium{}, \FiveCertAPHeavy{}: routing, not forgetting.

\paragraph{Routes from the example alone.}\label{sec:fromexample} With instructions, \recog{} reads
the example alone, so its routes are the same at every wording (Table~\ref{tab:instanceroutesshort}).
With the original wording \recog{} ends at \AudUThreeOriginalRidgeAPOne{} against \AudUThreeOriginalCertAPOne{}
for \cert{}, and \AudUThreeOriginalRidgeChanged{} and \AudUThreeOriginalHandicapChanged{} of \AudUThreeCompared{}
answers change (\recog{} and \handicap{}) where the zones change none. With heavy rewording \recog{} ends
at \AudUThreeHeavyRidgeAPOne{} and \handicap{} at \AudUThreeHeavyHandicapAPOne{}, against \AudUThreeHeavyCertAPOne{} for
\cert{}, with a ceiling of \AudUThreeHeavyOwnAPOne{}. By lines fixed before the runs it repairs the weak
spot of the zones. Three seeds, machine-written rewordings, and the benchmark gives the example
apart from the instruction (Appendix~\ref{sec:exampleroutes}).

\begin{table}[t]
  \caption{Five tasks with instructions, original and heavily reworded; re-test, means over three seeds (changed answers summed). Zones read the instruction, \recog{} and \handicap{} the example alone. Last rows: the published methods.}
  \label{tab:instanceroutesshort}
  \centering
  {\footnotesize\setlength{\tabcolsep}{3pt}\begin{tabular}{lrrrr}
    \toprule
    & \multicolumn{2}{c}{Final score $\uparrow$} & \multicolumn{2}{c}{\shortstack{Answers changed\\of \AudUThreeCompared{} $\downarrow$}} \\
    \cmidrule(lr){2-3}\cmidrule(l){4-5}
    & original & heavy & original & heavy \\
    \midrule
    \cert{} (zones on the instruction) & \AudUThreeOriginalCertAP{} & \AudUThreeHeavyCertAP{} & \AudUThreeOriginalCertChanged{} & \AudUThreeHeavyCertChanged{} \\
    \recog{} (example alone) & \AudUThreeOriginalRidgeAP{} & \AudUThreeHeavyRidgeAP{} & \AudUThreeOriginalRidgeChanged{} & \AudUThreeHeavyRidgeChanged{} \\
    \handicap{}, $\nu = 2$ (example alone) & \AudUThreeOriginalHandicapAP{} & \AudUThreeHeavyHandicapAP{} & \AudUThreeOriginalHandicapChanged{} & \AudUThreeHeavyHandicapChanged{} \\
    every input to its own task (ceiling) & \AudUThreeOriginalOwnAP{} & \AudUThreeHeavyOwnAP{} & \AudUThreeOriginalOwnChanged{} & \AudUThreeHeavyOwnChanged{} \\
    \midrule
    GainLoRA & \AudUThreeOriginalGainloraAP{} & \AudUThreeHeavyGainloraAP{} & \AudUThreeOriginalGainloraChanged{} & \AudUThreeHeavyGainloraChanged{} \\
    InfLoRA & \AudUThreeOriginalInfloraAP{} & \AudUThreeHeavyInfloraAP{} & \AudUThreeOriginalInfloraChanged{} & \AudUThreeHeavyInfloraChanged{} \\
    \bottomrule
\end{tabular}}

\end{table}

\paragraph{What the zones add, and what they cost.}\label{sec:stress}\label{sec:many} In a constructed case
where a sixth task arrives with SST-2 sentences and a rewording of task 2's definition, \nearest{} sends
all 100 old inputs with the light rewording to the new task, and nothing announces it; under \cert{} no
old input changes its route. The certified inputs of task 2
keep their route in every case (routes only; Appendix~\ref{app:fulltext}). The price is measured on the
\ManyTasks{} tasks of the default English split of Super-NaturalInstructions, in random order, with
zones only and no training. No certified input changed its route. But \ManyColliding{} tasks
(\ManyCollidingPct{}\%; \ManyCollidingPctLow{}\% to \ManyCollidingPctHigh{}\% in four other orders)
collided with an older task, and \ManyAllHeld{} were blocked on every input; most are families of
nearly identical definitions. On unedited instructions \nearest{} gives
\ManyNearRightPct{}\% of the inputs to their own task, \cert{} \ManyCertRightPct{}\%. A variant that serves
the siblings, \exactfirst{}, was designed after this result, and serves all but \ExactTasksNone{} tasks
(not an independent test). On machine paraphrases of \WidenTasksEval{} tasks the zones certify only
\WidenRefTfivePLow{}\% to \WidenRefTfivePHigh{}\% of inputs for their own task, none of five attempts to
widen them met the bar set for it, and \nearest{} changed \RerouteNearNChanged{} of \RerouteInputs{} routes
after all tasks had arrived (we had expected \RerouteExpectNearLow{}\% to \RerouteExpectNearHigh{}\%),
the same as \cert{}. A plain nearest signature gives \RerouteSigOwnEnd{}\% of these inputs to their own task,
\cert{} \RerouteCertOwnEnd{}\%; a fallback found at the end, scaled by the distance to unrelated text,
gives \RerouteCertFourItOwn{}\% behind the zones and \RerouteBgItOwn{}\% alone (fourth pivot language; routes only). There the zones add the proof, not better routing (Appendices~\ref{app:many},
\ref{app:widen} and~\ref{app:reroute}).\label{sec:edits}\label{sec:models} Large mechanical edits of the instruction lose the
certificate (with sixteen typos only 45.4\% of the inputs reach their own task; Appendix~\ref{app:edits});
routes with other models are in Appendix~\ref{app:models}.
